\section*{Capítulo \ref{ch:b2:angiosperm-reproduction} --- Reproducción sexual de las plantas con flores}

\begin{solution}{exo:b2:angiosperm-reproduction:1}
Sépalos (cáliz), pétalos (corola), \omterm{def:b2:angiosperm-reproduction:flower}{estambres} (androceo), \omterm{def:b2:angiosperm-reproduction:flower}{carpelos}
(gineceo). \omterm{def:b2:angiosperm-reproduction:flower}{Estambre}: filamento y antera. \omterm{def:b2:angiosperm-reproduction:flower}{Carpelo}: \omterm{def:b2:angiosperm-reproduction:flower}{estigma}, estilo y
\omterm{def:b2:angiosperm-reproduction:flower}{ovario} con óvulos. Todas estas partes son \omterm{def:b2:plant-life-cycles:cycles}{esporofito} \omterm{def:b2:meiosis-heredity:meiosis}{diploide}; solo el
\omterm{prop:b2:angiosperm-reproduction:pollen}{grano de polen} (dentro de la antera) y el \omterm{prop:b2:angiosperm-reproduction:embryosac}{saco embrionario} (dentro del
óvulo) son \omterm{def:b2:plant-life-cycles:cycles}{gametofito} \omterm{def:b2:meiosis-heredity:meiosis}{haploide}.
\end{solution}

\begin{solution}{exo:b2:angiosperm-reproduction:2}
\omterm{def:b2:plant-life-cycles:heterospory}{Microspora} $n$; célula vegetativa $n$; célula espermática $n$; oosfera
$n$; núcleo polar $n$; cigoto $2n$; endospermo $3n$; \omterm{def:b2:angiosperm-reproduction:seed}{cubierta seminal}
$2n$ (tegumentos maternos); pared del \omterm{def:b2:angiosperm-reproduction:seed}{fruto} $2n$ (\omterm{def:b2:angiosperm-reproduction:flower}{ovario} materno).
\end{solution}

\begin{solution}{exo:b2:angiosperm-reproduction:3}
Viento (gramíneas, roble, abedul, avellano, llantén): \omterm{def:b2:angiosperm-reproduction:flower}{flores} pequeñas y
verdes, sin pétalos, ni olor, ni \omterm{def:b2:angiosperm-reproduction:pollination}{néctar}, anteras colgantes sobre largos
filamentos, \omterm{def:b2:angiosperm-reproduction:flower}{estigmas} plumosos, enormes cantidades de \omterm{def:b2:plant-life-cycles:heterospory}{polen} liso y seco,
floración antes de las hojas. Abeja (trébol, salvia, dedalera, lavanda,
manzano): pétalos coloreados, a menudo con guías ultravioletas,
plataforma de aterrizaje, perfume, \omterm{def:b2:angiosperm-reproduction:pollination}{néctar}, \omterm{def:b2:plant-life-cycles:heterospory}{polen} pegajoso y esculpido en
cantidad moderada, floración de día.
\end{solution}

\begin{solution}{exo:b2:angiosperm-reproduction:4}
Las dos células espermáticas de un mismo \omterm{thm:b2:angiosperm-reproduction:tube}{tubo polínico} se fusionan, una
con la oosfera (cigoto, $2n$, el embrión) y otra con la célula central
(núcleo del endospermo, $3n$, la reserva de alimento). El endospermo solo
empieza a crecer tras la fecundación, de modo que la planta solo
aprovisiona los óvulos que contienen un embrión, a diferencia de la
gimnosperma, que construye antes su reserva de alimento.
\end{solution}

\begin{solution}{exo:b2:angiosperm-reproduction:5}
$25/1.2 = \qty{21}{h}$. Las sedas que emergen los días 5, 6 y 7 encuentran
\omterm{def:b2:plant-life-cycles:heterospory}{polen}; las de los días 8 a 12, no: se poliniza $3/8$ de las sedas, y la
mazorca queda estéril en cinco octavas partes.
\end{solution}

\begin{solution}{exo:b2:angiosperm-reproduction:6}
$S_1S_2\times S_1S_2$: 0. $S_1S_2\times S_1S_3$: $\tfrac12$ (solo crece
el \omterm{def:b2:plant-life-cycles:heterospory}{polen} $S_3$). $S_1S_2\times S_3S_4$: todo. Descendientes del segundo
cruzamiento: oosferas $S_1$ o $S_2$, células espermáticas $S_3$: $S_1S_3$
y $S_2S_3$, la mitad de cada uno.
\end{solution}

\begin{solution}{exo:b2:angiosperm-reproduction:7}
Una planta lleva 2 de los $n$ \omterm{def:b2:meiosis-heredity:vocabulary}{alelos} y rechaza el \omterm{def:b2:plant-life-cycles:heterospory}{polen} que lleve
cualquiera de ellos: fracción compatible $(n - 2)/n$. Ningún pistilo
rechaza un \omterm{def:b2:meiosis-heredity:vocabulary}{alelo} nuevo, de modo que su \omterm{def:b2:plant-life-cycles:heterospory}{polen} engendra más \omterm{def:b2:angiosperm-reproduction:seed}{semillas} que la
media y el \omterm{def:b2:meiosis-heredity:vocabulary}{alelo} se propaga hasta ser tan común como los demás; la misma
ventaja protege de la pérdida a todo \omterm{def:b2:meiosis-heredity:vocabulary}{alelo} raro. La selección acumula
por tanto \omterm{def:b2:meiosis-heredity:vocabulary}{alelos}, y las poblaciones naturales de especies
autoincompatibles llevan decenas.
\end{solution}

\begin{solution}{exo:b2:angiosperm-reproduction:8}
Los dos \omterm{prop:b2:angiosperm-reproduction:embryosac}{núcleos polares} son copias de una misma \omterm{def:b2:plant-life-cycles:heterospory}{megaspora}, de modo que la
célula central es $YY$ o $yy$ (la mitad de cada una); la célula
espermática es $Y$ o $y$ (la mitad de cada una): endospermo $YYY$, $YYy$,
$Yyy$, $yyy$, cada uno con $\tfrac14$; tres cuartas partes amarillos.
Polinizada por $yy$: $YYy$ e $yyy$, la mitad de cada uno --- la mitad de
los granos amarillos.
\end{solution}

\begin{solution}{exo:b2:angiosperm-reproduction:9}
Se cortan granos de cebada por la mitad; las mitades con embrión digieren
su almidón y las mitades sin embrión no, salvo que se añada giberelina,
en cuyo caso sí lo hacen. (a) El embrión es la fuente de la señal, no de
la enzima. (b) La giberelina es la señal, suficiente por sí sola. (c) La
amilasa la fabrica la capa de aleurona del endospermo, que responde a la
hormona. El medio grano sin embrión separa la señal de la respuesta: sin
él no podría saberse si el embrión digiere el almidón por sí mismo.
\end{solution}

\begin{solution}{exo:b2:angiosperm-reproduction:10}
Viento: $10^{6}\times\qty{1}{\micro g} = \qty{1}{g}$ por óvulo. Animal:
\qty{1}{mg} de \omterm{def:b2:plant-life-cycles:heterospory}{polen} más \qty{0.5}{mg} de azúcar $= \qty{1.5}{mg}$
por óvulo --- unas setecientas veces más barato. La \omterm{def:b2:angiosperm-reproduction:pollination}{polinización} por el
viento no necesita pareja: funciona a principios de primavera, antes de
que vuelen los insectos, en lugares fríos, ventosos y abiertos, de noche
y bajo la lluvia, en rodales densos de una sola especie, y no pueden
engañarla los ladrones de \omterm{def:b2:angiosperm-reproduction:pollination}{néctar} ni abandonarla un polinizador que se
extingue.
\end{solution}

\begin{solution}{exo:b2:angiosperm-reproduction:11}
Las \omterm{def:b2:angiosperm-reproduction:seed}{semillas} que caen al pie del progenitor quedan a su sombra, entre sus
raíces y donde se concentran sus depredadores de \omterm{def:b2:angiosperm-reproduction:seed}{semillas} y sus
patógenos específicos; casi todas mueren. Un ave lleva una \omterm{def:b2:angiosperm-reproduction:seed}{semilla} a
kilómetros (\qty{12}{km} en veinte minutos) y la deja caer, con abono, en
un seto o un claro; aunque solo una \omterm{def:b2:angiosperm-reproduction:seed}{semilla} de cada cien caiga en un buen
sitio, eso supera a todas las \omterm{def:b2:angiosperm-reproduction:seed}{semillas} que quedan bajo el árbol, y la
población se extiende a la velocidad de las aves, no a la de la fruta que
cae. La pulpa es el precio del billete.
\end{solution}

\begin{solution}{exo:b2:angiosperm-reproduction:12}
El encierro obliga al \omterm{def:b2:plant-life-cycles:heterospory}{polen} a crecer a través de tejido materno, lo que
permite al pistilo ponerlo a prueba (autoincompatibilidad, rechazo de
especies ajenas), dejar que muchos tubos compitan para que el más rápido
engendre la \omterm{def:b2:angiosperm-reproduction:seed}{semilla}, proteger los óvulos de la desecación y de los
herbívoros, y convertir el \omterm{def:b2:angiosperm-reproduction:flower}{ovario} en un \omterm{def:b2:angiosperm-reproduction:seed}{fruto} que recluta animales para
la dispersión; la \omterm{prop:b2:angiosperm-reproduction:double}{doble fecundación} aprovisiona después solo los óvulos
fecundados. Los costes: un \omterm{prop:b2:angiosperm-reproduction:pollen}{grano de polen} debe llegar a un \omterm{def:b2:angiosperm-reproduction:flower}{estigma} y no
al propio óvulo, de modo que hacen falta un estilo, un tubo y su guiado,
y el \omterm{def:b2:angiosperm-reproduction:seed}{fruto} es una inversión cuantiosa. El balance favoreció tanto a las
angiospermas que pasaron de no existir a representar nueve décimas partes
de las especies vegetales en cien millones de años.
\end{solution}

\begin{solution}{pb:b2:angiosperm-reproduction:1}
\textbf{1.} Sobre A ($S_1S_2$): \omterm{def:b2:plant-life-cycles:heterospory}{polen} de A $f = 0$; de B ($S_1$, $S_3$) $f =
\tfrac12$; de C ($S_3$, $S_4$) $f = 1$.
\textbf{2.} Sobre B ($S_1S_3$): A $\tfrac12$, B 0, C $\tfrac12$. Sobre C
($S_3S_4$): A 1, B $\tfrac12$, C 0.
\textbf{3.} 50 granos compatibles; $0.1\times 50 = 5$ óvulos.
\textbf{4.} Desde C: 100 compatibles, $0.1\times 100 = 10$, es decir, los
diez óvulos; desde A: ninguno.
\textbf{5.} Se conservan los \omterm{def:b2:angiosperm-reproduction:seed}{frutos} de las visitas desde C (10 \omterm{def:b2:angiosperm-reproduction:seed}{semillas});
los de B (5 \omterm{def:b2:angiosperm-reproduction:seed}{semillas}), por los pelos; los de A caen. Un huerto solo de
A no da nada.
\textbf{6.} Cuaja la mitad de las \omterm{def:b2:angiosperm-reproduction:flower}{flores} (las visitadas desde C): 2500
\omterm{def:b2:angiosperm-reproduction:seed}{frutos} por árbol, diez veces los 250 necesarios --- una cosecha completa,
y el árbol dejará caer el excedente.
\textbf{7.} Libera \omterm{def:b2:plant-life-cycles:heterospory}{polen} compatible durante toda la floración de todas
las variedades; pero si comparte sus dos \omterm{def:b2:meiosis-heredity:vocabulary}{alelos} $S$ con una variedad, no
poliniza nada en ella.
\textbf{8.} $8\times 10^{7}/(7\times 86400) = 132$ granos por metro
cuadrado y por segundo.
\textbf{9.} $132/0.25 = 530$ granos por metro cúbico.
\textbf{10.} $132\times 10^{-4} = 0.013$ por segundo: un grano cada
\qty{76}{s}.
\textbf{11.} $0.013\times 86400 \approx 1100$ granos al día; el primer tubo
que llega al óvulo lo fecunda, y el óvulo queda entonces cerrado a los
demás.
\textbf{12.} \qty{20}{h} después de la \omterm{def:b2:angiosperm-reproduction:pollination}{polinización}.
\textbf{13.} \omterm{def:b2:plant-life-cycles:heterospory}{Polen} los días 0--7, sedas a partir del día 5: solo se
solapan los días 5, 6 y 7 --- $3/8$ de las sedas, y las
sedas que emergen después del día 7 no reciben nada; la mazorca queda en
gran parte estéril. La sequía en la floración femenina es la causa
clásica de una cosecha de maíz fallida.
\textbf{14.} Embrión $Yy$; endospermo $YYy$.
\textbf{15.} Amarillo, porque $Y$ es \omterm{def:b2:meiosis-heredity:vocabulary}{dominante}: el carácter del progenitor
que aporta el \omterm{def:b2:plant-life-cycles:heterospory}{polen} se manifiesta en la \omterm{def:b2:angiosperm-reproduction:seed}{semilla} sobre la planta madre ---
la xenia.
\textbf{16.} Célula central $YY$ o $yy$: endospermo $YYy$ (la mitad) e
$yyy$ (la mitad).
\textbf{17.} $YYy$: 20 unidades; $Yyy$: 10 unidades; un $yyY$ con la $Y$
procedente del padre es el mismo \omterm{def:b2:meiosis-heredity:vocabulary}{genotipo} $Yyy$, 10 unidades. Dos dosis
paternas no pueden existir: una sola célula espermática fecunda la célula
central, y el endospermo tiene siempre dos genomas maternos por uno
paterno.
\textbf{18.} Es la reserva de almidón; $0.3/0.33 = \qty{91}{\%}$ del
grano.
\textbf{19.} La reserva solo se construye tras la fecundación, de modo
que no se destinan recursos a óvulos sin embrión; el \omterm{def:b2:plant-life-cycles:cycles}{gametofito} de la
gimnosperma se construye antes y se desperdicia si la \omterm{def:b2:angiosperm-reproduction:pollination}{polinización} falla.
\textbf{20.} $0.9\times 600 = 540$ granos por mazorca y por planta;
$4320$ por metro cuadrado; $4320\times 0.3 = \qty{1.3}{kg}$ de endospermo
por metro cuadrado.
\textbf{21.} \qty{13}{t/ha}.
\textbf{22.} Carbono del grano $0.5\times 1.2 = \qty{0.6}{kg/m^2}$;
carbono del endospermo $0.45\times 1.3 = \qty{0.58}{kg/m^2}$: coherente
(el pequeño resto corresponde al embrión y la cubierta).
\textbf{23.} $8\times 10^{7}/4320 \approx \num{18500}$ \omterm{prop:b2:angiosperm-reproduction:pollen}{granos de polen}
por grano de maíz.
\textbf{24.} Su propia nube de \omterm{def:b2:plant-life-cycles:heterospory}{polen} es tenue y se la lleva el viento, y
su panícula libera la mayor parte del \omterm{def:b2:plant-life-cycles:heterospory}{polen} antes de que emerjan sus
propias sedas, de modo que muchas sedas nunca se polinizan y la mazorca
queda con huecos.
\textbf{25.} Bloque A: cuaja la mitad de las \omterm{def:b2:angiosperm-reproduction:flower}{flores}, 2500 \omterm{def:b2:angiosperm-reproduction:seed}{frutos} por
árbol, una cosecha completa. Campo: 4320 granos por metro cuadrado,
\qty{13}{t/ha}.
\end{solution}
